\documentclass[11pt]{article}
\usepackage[margin=1in]{geometry}
\usepackage{amsmath,amssymb,amsthm,mathtools}
\usepackage[hidelinks]{hyperref}
\newtheorem{theorem}{Theorem}
\newtheorem{lemma}[theorem]{Lemma}
\newcommand{\spen}{s^{+}}
\title{Positive Square Energy of Theta Graphs}
\author{math-god}
\date{23 July 2026}
\begin{document}
\maketitle

\begin{abstract}
For a graph $G$ with adjacency eigenvalues $\lambda_i$, let
$s^+(G)=\sum_{\lambda_i>0}\lambda_i^2$.  We prove that every simple theta
graph $G$ satisfies
$s^+(G)>|V(G)|$.  This verifies the theta-graph case of the conjecture
of Akbari, Kumar, Mohar, Pragada, and Zhang~\cite{AKMPZ} that every connected $n$-vertex
graph with at least $n+1$ edges has $s^+\ge n$.  The proof combines bipartite
spectral symmetry, an improved induced-$P_3$ removal lemma, and explicit
positive-semidefinite variational witnesses for the three short exceptional
parity families.
\end{abstract}

\section{Introduction}
A simple theta graph $\Theta(a,b,c)$ consists of three internally
vertex-disjoint paths of positive integer edge lengths $a,b,c$ between the same pair of branch
vertices; at most one of $a,b,c$ is one.  It has
\[
 n=a+b+c-1,\qquad m=a+b+c=n+1.
\]
Our result is the following.

\begin{theorem}\label{thm:main}
Every simple theta graph $G$ satisfies $s^+(G)>|V(G)|$.
\end{theorem}

We use two standard facts.  If $G$ is bipartite, its adjacency spectrum is
symmetric; together with $\operatorname{tr}(A(G)^2)=2|E(G)|$, this gives
\begin{equation}\label{eq:bip}
 s^+(G)=|E(G)|.
\end{equation}
For a real symmetric matrix $M$, define $s^+(M)$ analogously.  Then
\begin{equation}\label{eq:var}
 s^+(M)=\max_{X\succeq0}\{2\operatorname{tr}(MX)-\operatorname{tr}(X^2)\}.
\end{equation}
The latter follows by diagonalizing $M$.

We require the following rigorous form of the improved $P_3$-removal lemma of
Akbari et al.~\cite[Lemma~2.4]{AKMPZ}.  Their proof reduces the improvement to a scalar inequality
checked numerically; we include an exact verification.

\begin{lemma}[Improved $P_3$ removal]\label{lem:p3}
If three vertices $U$ of a graph $G$ induce a $P_3$, then some $v\in U$
satisfies
\begin{equation}\label{eq:p3}
 s^+(G)\ge s^+(G-v)+\frac{17}{16}.
\end{equation}
\end{lemma}

\begin{proof}
Put $C=17/16$ and
\[
B=\begin{pmatrix}0&1&0\\1&0&1\\0&1&0\end{pmatrix}.
\]
For a symmetric matrix $Z$, let
$r_i(Z)=Z_{ii}^2+2\sum_{j\ne i}Z_{ij}^2$, the Frobenius contribution removed
when row and column $i$ are deleted.  We show that for every
$3\times3$ matrix $M\succeq0$, some $i$ has $r_i(B+M)>C$.

Conjugating by $\operatorname{diag}(1,-1,1)$ reduces this to $B-M$.  Write
\[
M=\begin{pmatrix}a&b&d\\b&c&e\\d&e&f\end{pmatrix}\succeq0
\]
and suppose all three row/column sums are at most $C$.  Explicitly,
\begin{align*}
a^2+2(1-b)^2+2d^2&\le C,\\
c^2+2(1-b)^2+2(1-e)^2&\le C,\\
f^2+2d^2+2(1-e)^2&\le C.
\end{align*}
With $x=(b+e)/2$, convexity and Cauchy--Schwarz give
\begin{align}
c^2&\le C-4(1-x)^2,\label{eq:p3a}\\
(a+f+2d)^2&\le6\bigl(C-2(1-x)^2\bigr).\label{eq:p3b}
\end{align}
Moreover $M\succeq0$ gives
\begin{equation}\label{eq:p3c}
c(a+f+2d)\ge(b+e)^2=4x^2.
\end{equation}
The initial inequalities imply $15/32\le x\le17/16$: the lower bound follows
from $|1-x|\le\sqrt C/2$ and $\sqrt C\le C$, while $b^2\le ac$,
$e^2\le cf$, and $a,c,f\le\sqrt C$ give $x\le C$.
Equations \eqref{eq:p3a}--\eqref{eq:p3c} would imply
\[
16x^4\le6\bigl(C-4(1-x)^2\bigr)\bigl(C-2(1-x)^2\bigr).
\]
The reverse inequality holds exactly.  Its left side minus right side is
\[
F(x)=-32x^4+192x^3-\frac{999}{4}x^2
+\frac{231}{2}x-\frac{2115}{128}.
\]
Under affine parametrization of $[15/32,2/3]$, its degree-four Bernstein
coefficients are
\[
\frac{31815}{32768},\frac{5821}{8192},\frac{2515}{8192},
\frac{37}{2304},\frac{461}{10368};
\]
on $[2/3,17/16]$ they are
\[
\frac{461}{10368},\frac{175}{1728},\frac{26459}{18432},
\frac{35027}{6144},\frac{28193}{2048}.
\]
All are positive, while Bernstein basis functions are nonnegative.  This
proves the local matrix assertion.

Finally,
\[
s^+(H)=\min_{R\succeq0}\|A(H)+R\|_F^2.
\]
Take a minimizer for $G$ and restrict it to the principal block on $U$.  The
local assertion supplies a vertex whose full row/column contribution exceeds
$17/16$; deleting it and using the remaining principal submatrix as a
competitor for $G-v$ proves \eqref{eq:p3}.
\end{proof}

\section{The parity reduction}
A theta graph is bipartite exactly when $a,b,c$ have the same parity; then
\eqref{eq:bip} gives $s^+=m=n+1$.

Suppose it is nonbipartite.  Relabel so that $a\equiv b\pmod2$ and $c$ has the
opposite parity.  If $c\ge4$, choose three consecutive internal vertices on
the $c$-path.  They induce a $P_3$.  Deleting any of the three breaks the
$c$-path and leaves a connected bipartite unicyclic graph: its unique cycle
has length $a+b$, which is even.  Every possible deletion has $n-1$ edges, so
\eqref{eq:bip} and Lemma~\ref{lem:p3} give
\begin{equation}\label{eq:long}
 s^+(G)\ge n-1+\frac{17}{16}=n+\frac1{16}.
\end{equation}
It remains to treat
\[
 \Theta(2r,2s,1)\ (r,s\ge1),\quad
 \Theta(2r+1,2s+1,2)\ (r,s\ge0,\ r+s\ge1),
\]
and $\Theta(2r,2s,3)$ with $r,s\ge1$.

\section{A bordered-cycle witness}
We first settle the middle family.  Let $A$ be the adjacency matrix of an
even cycle $C_N$, let $u,v$ lie in different bipartition classes, and put
$b=e_u+e_v$.  Adding a common neighbor gives
\[
 M=\begin{pmatrix}A&b\\b^T&0\end{pmatrix}.
\]
Write $P=A_+$ for the positive spectral part, and put $p=b^TPb$ and
$q=b^TP^2b$.  The matrix
\[
 X=\begin{pmatrix}P&Pb/2\\b^TP/2&p/4\end{pmatrix}\succeq0
\]
has a Gram factorization with block rows $P^{1/2}$ and
$b^TP^{1/2}/2$.  Substitution in \eqref{eq:var} gives
\[
 s^+(M)\ge N+2p-\frac q2-\frac{p^2}{16}
 \ge N+p-\frac{p^2}{16},
\]
because $0\preceq P\preceq2I$ implies $q\le2p$.

For $N\ge6$, the second and fourth cycle spectral moments are $2$ and $6$;
the latter counts the six closed four-step walks from a cycle vertex.
H\"older interpolation gives
\[
 \frac1N\operatorname{tr}|A|\ge\frac2{\sqrt3}.
\]
Since $u,v$ are in opposite classes, $|A|_{uv}=0$, and
\[
p=\frac1N\operatorname{tr}|A|+\mathbf1_{\{u\sim v\}}\ge\frac2{\sqrt3}.
\]
Cauchy--Schwarz gives $N^{-1}\operatorname{tr}|A|\le\sqrt2$, so $p<3$.
As $x-x^2/16$ increases on $[0,3]$,
\[
 p-\frac{p^2}{16}\ge\frac2{\sqrt3}-\frac1{12}>1.
\]
In $C_4$ every opposite-class pair is adjacent, so $p=2$ and the estimate is
stronger.  Therefore
\begin{equation}\label{eq:oddodd2}
 s^+(\Theta(2r+1,2s+1,2))>|V|.
\end{equation}

\section{The chord family}
Write $\Theta(2r,2s,1)$ as an even cycle $C_N$ plus a chord between
same-class vertices $u,v$.  Put
\[
 a=\frac{e_u+e_v}{\sqrt2},\quad b=\frac{e_u-e_v}{\sqrt2},\quad
 E=aa^T-bb^T,
\]
and, for $P=A_+$,
\[
 x=a^TPa,\quad y=b^TPb,\quad q=b^TP^2b.
\]
The dihedral reflection interchanging $u,v$ commutes with $P$, fixes $a$, and
negates $b$; hence $a^TPb=0$.  With
\[
 T=I-\frac13bb^T,\qquad \alpha=(1-x)_+,
\]
the matrix $X=TPT+\alpha aa^T$ is positive semidefinite.  Exact expansion in
\eqref{eq:var} yields
\begin{equation}\label{eq:chordgain}
 s^+(A+E)-N\ge
 2x-\frac89y-\frac29q-\frac{25}{81}y^2+(1-x)_+^2.
\end{equation}

Bipartite symmetry and spectral Cauchy--Schwarz give $y^2\le q/2$.  Indeed,
if $\Pi_+$ is the positive spectral projector, then
\[
y^2\le(b^T\Pi_+b)(b^TP^2b),\qquad b^T\Pi_+b\le\frac12,
\]
because the bipartition involution exchanges positive and negative spectral
subspaces.  On one bipartition class $P^2$ has quadratic form $A^2/2$, so
\[
q=\frac12b^TA^2b=1-\frac{(A^2)_{uv}}2\le1.
\]
Finally $a^TAa=0$ and $|A|\succeq A^2/2$ give
\[
x=\frac12a^T|A|a\ge\frac14a^TA^2a
=\frac{2+(A^2)_{uv}}4\ge\frac12.
\]
Thus $y\le1/\sqrt2$, $y^2\le1/2$, and
$2x+(1-x)_+^2\ge5/4$.  Hence the right side of
\eqref{eq:chordgain} is at least
\[
 \frac54-\frac8{9\sqrt2}-\frac29-\frac{25}{162}
 =\frac{283}{324}-\frac{4\sqrt2}{9}>0.
\]
Thus
\begin{equation}\label{eq:eveneven1}
 s^+(\Theta(2r,2s,1))>|V|.
\end{equation}

\section{The length-three ear}
Let $A_0=A(C_N)\oplus A(K_2)$, and let $B$ be the symmetric adjacency matrix
of the two matching edges joining the endpoints of $K_2$ to same-class cycle
vertices $u,v$.  The theta adjacency matrix is $A=A_0+B$.
Put $X_0=(A_0)_+$ and $S=\operatorname{tr}(X_0^2)=N+1$.  For $t\ge0$ let
\[
 X_t=(I+tB)X_0(I+tB)\succeq0.
\]
Scaling a fixed witness in \eqref{eq:var} gives, provided
$\operatorname{tr}(AX_t)>0$,
\begin{equation}\label{eq:ratio}
 s^+(A)\ge\frac{\operatorname{tr}(AX_t)^2}{\operatorname{tr}(X_t^2)}.
\end{equation}
By vertex transitivity write
$a=(A(C_N)_+)_{uu}=(A(C_N)_+)_{vv}$ and
$b=(A(C_N)_+)_{uv}$.  Block multiplication gives
\begin{align}
 \operatorname{tr}(AX_t)&=S+(2+4a)t+2bt^2,\label{eq:num}\\
 \operatorname{tr}(X_t^2)&=S+(6+8a+8b)t^2
 +(1+2a^2+2b^2)t^4.\label{eq:den}
\end{align}
We have $1/2\le a\le1/\sqrt2$ and $|b|\le a$: use
$a=|A(C_N)|_{uu}/2$, the inequalities
$A(C_N)^2/2\preceq|A(C_N)|$ and
$|A(C_N)|_{uu}\le\sqrt{(A(C_N)^2)_{uu}}=\sqrt2$, and the
$\{u,v\}$ principal minor of $P\succeq0$.

Take $t=1/3$; positivity of $\operatorname{tr}(AX_t)$ follows from
\eqref{eq:num} and $|b|\le a$.  Let $U,D$ denote the quantities in
\eqref{eq:num}--\eqref{eq:den}.  Direct expansion of
$E=U^2-(S+1)D$ gives
\begin{align*}
81E={}&-2Sa^2+144Sa-2Sb^2-36Sb-28S\\
&+142a^2+48ab+72a+2b^2-48b-19.
\end{align*}
The coefficient of $S$ is at least $-4a^2+108a-28\ge25$, so this increases
with $S\ge5$.  At $S=5$ it is at least
\[
132a^2+48ab+792a-8b^2-228b-159.
\]
Its derivative in $b$ is $48a-16b-228<0$, so it decreases on $[-a,a]$.
At $b=a$ it is
$172a^2+564a-159\ge166$.  Thus $E>0$, and \eqref{eq:ratio} proves
\begin{equation}\label{eq:eveneven3}
 s^+(\Theta(2r,2s,3))>N+2=|V|.
\end{equation}

Equations \eqref{eq:long}, \eqref{eq:oddodd2},
\eqref{eq:eveneven1}, and \eqref{eq:eveneven3}, together with the bipartite
case, prove Theorem~\ref{thm:main}.

\section{Verification}
The source repository contains the complete attack notebook, detailed proof
notes, and the exact script \texttt{experiments/theta\_paper\_certificate.py}
checking the Bernstein coefficients and displayed trace identities.  These
computations are supplementary; the proofs above do not depend on finite
numerical verification.

\begin{thebibliography}{9}
\bibitem{AKMPZ}
S.~Akbari, H.~Kumar, B.~Mohar, S.~Pragada, and S.~Zhang,
\emph{Refinement of a conjecture on positive square energy of graphs},
arXiv:2506.07264v1, 2025.
\bibitem{EFGW}
C.~Elphick, M.~Farber, F.~Goldberg, and P.~Wocjan,
\emph{Conjectured bounds for the sum of squares of positive eigenvalues of a
graph}, Discrete Math. 339 (2016), 2215--2223.
\bibitem{AbiadEtAl}
A.~Abiad, L.~de Lima, D.~N.~Desai, K.~Guo, L.~Hogben, and J.~Madrid,
\emph{Positive and negative square energies of graphs}, Electron. J. Linear
Algebra 39 (2023), 307--326.
\end{thebibliography}
\end{document}
